\chapter{Memory as a Limiting Case: MEM|8 and the Ontology of Traces}
\label{ch:mem8}

\begin{quotation}
\noindent\textit{Synopsis.} This chapter treats memory as both an example
and a challenge to the framework. Two readings of MEM|8 are set out: a
narrow reading, under which memory is simply another reconstructive process
operating over traces, and a broad reading, under which memory or process
is ontologically prior to representation itself, implying that all access
to artifacts is already reconstructive. The chapter does not attempt to
resolve this dispute. It instead analyzes the consequences of each reading
and identifies which parts of the framework remain stable regardless of the
outcome, treating the fork as an explicit, acknowledged dependency rather
than concealing it behind a premature resolution.
\end{quotation}

This chapter is a friction chapter, not a worked example, and is written
to be read as such. It does not resolve the question it raises; it states
the question precisely enough that the rest of the monograph's dependence
on one answer or the other is visible, and it closes with a table stating
exactly which results survive regardless of the answer and which do not.

\section{Two Readings}

MEM|8's core claim --- that memory is prior to representation rather than
a degraded copy of it --- admits at least two readings with very
different consequences for this monograph, and the difference between
them is worth setting out in full before either is evaluated.

Under the narrow reading, biological recall specifically is
reconstructive: a trace-based regeneration process rather than retrieval
of a stored copy. This reading is well supported by the eyewitness-memory
literature and is compatible with treating memory as one more point on
the fidelity continuum of Chapter~\ref{ch:fidelity-continuum}, alongside
photocopies, scans, and decoded audio. On this reading, memory is unusual
among renderings only in degree --- its fidelity tends to be low and its
schema-driven gap-filling tends to be confident about its own accuracy in
a way a photocopier is not --- not in kind.

Under the broad reading, representation \emph{in general} is a
stabilization of prior memory or process, not the reverse: perception,
recognition, and even the sense of directly encountering an external
object are themselves outputs of an ongoing memory-like process, rather
than a separate, more basic capacity that memory merely draws on
afterward. Under this reading, no consumer --- including a chess player
looking at a physical board --- ever has strictly unmediated access to an
external artifact, since the very act of perceiving the board is already
a memory-like construction rather than a transparent window onto it. This
would undermine the very idea of a fixed root object
(\cref{def:root-object}) for any embodied consumer, since a root object's
usefulness (\cref{prop:root-recovery}) depends on there being a stable
thing to re-derive from, and the broad reading denies that anything
external is ever encountered other than through the same reconstructive
process that governs recall.

\section{Considerations on Each Side}

It would be premature, and beyond this monograph's scope, to adjudicate
between these readings, but it is worth surveying briefly what would need
to be true for each to hold, so that the fork is a substantive question
rather than an empty one. The narrow reading is the more conservative
position and requires only what is already well established:
reconstructive-memory research shows that recall is schema-driven and
error-prone in characteristic, well-documented ways, without requiring
any claim about perception or representation more broadly. The broad
reading is a stronger, more general claim, and its most natural allies
are views in cognitive science that treat perception itself as a
constructive or predictive process rather than a passive registration of
external stimuli --- the general family of ideas sometimes described as
treating perception as a controlled or constrained hallucination, or as
governed by predictive processing in which incoming signal is used to
correct an ongoing internally generated model rather than read off
directly. If some view in this family is correct, the boundary between
``perceiving the board'' and ``remembering the board'' may not be a
boundary between two different kinds of access at all, but two points on
one continuous process --- which is exactly what the broad reading
claims and what this monograph's own Part~\ref{part:modes-of-access}
argued for the boundary between preservation and reconstitution more
narrowly (\cref{prop:preservation-limit}). This is a genuine structural
echo, and it is tempting to treat it as evidence for the broad reading;
this monograph resists that temptation, since an echo between two
frameworks is not the same as a proof that one licenses the other, and
the question of whether perception and memory share a single underlying
mechanism is an empirical matter this book is not positioned to settle.

\section{Consequences of the Broad Reading}

If the broad reading is correct, the admissibility and distinguishability
apparatus of Part~\ref{part:foundations} would need to be built on top of
a process-first ontology, rather than assuming a stable artifact that a
consumer either does or does not have access to. This monograph
explicitly declines to resolve this question here. It is stated as the
single largest open dependency of the entire framework.

\begin{openquestion}
\label{oq:mem8-fork}
Does MEM|8's claim generalize beyond biological recall (the narrow
reading) to representation as such (the broad reading)? If the broad
reading holds, Part~\ref{part:foundations}'s treatment of artifacts as
stable objects available for direct or indirect access requires revision
from the ground up, rather than the more modest correction made in
\cref{prop:preservation-limit}.
\end{openquestion}

\section{What Survives Regardless of the Fork}

Both readings agree that a trace or a percept is a finite, lossy
substrate; they disagree only about whether any \emph{non}-trace
substrate exists at all for an embodied consumer. This distinction is
sharp enough to sort the monograph's results cleanly into those that
depend on it and those that do not.

\begin{center}
\small
\begin{tabular}{p{4.3cm}p{2.3cm}p{2.3cm}}
\toprule
Result & Narrow reading & Broad reading \\
\midrule
\cref{prop:recursive-insufficiency} & holds & holds \\
\cref{thm:composition-failure} & holds & holds \\
\cref{prop:verification-sufficiency} (proofs) & holds & holds \\
\cref{prop:preservation-limit} & holds as stated & upper limit never reached \\
\cref{oq:unmediated-access} & resolved: no & open in stronger form \\
Ch.~\ref{ch:chess}'s board case & fidelity $\to 1$ available & fidelity $\to 1$ unavailable \\
Embodied cases, Part~\ref{part:worked-domains} & conditional & structurally revised \\
\bottomrule
\end{tabular}
\end{center}

The pattern in this table is not accidental. Results stated about
\emph{formal or machine} substrates --- a proof checked by an algorithm
against explicit axioms, a substrate chain defined over named
distinctions without reference to how those distinctions come to be
perceived --- do not mention embodied access at all, and so cannot be
disturbed by a dispute about what embodied access amounts to.
\cref{thm:composition-failure} in particular is stated entirely in terms
of renderings, budgets, and discard sets; nothing in its proof
(Appendix~\ref{app:proof}) invokes a consumer's perceptual access to
anything, so it is simply not the kind of claim the MEM|8 fork bears on.
Results that \emph{do} mention embodied access --- most sharply,
\cref{oq:unmediated-access} and every worked-domain discussion of a
physical object sitting in front of an observer --- are exactly the ones
the fork can move. The table's purpose is to make that boundary explicit
rather than leave it to be discovered piecemeal across later chapters.

\begin{remark}[The Fork's Scope, Precisely]
\label{rem:fork-scope}
A further distinction is worth making explicit, since it is easy to
conflate in practice and doing so caused an error corrected during the
drafting of this monograph: the fork bears on whether a consumer's
perceptual access to a \emph{given} rendering can reach fidelity $1$
relative to that rendering. It does not bear on whether that rendering
was already low-fidelity for independent, domain-structural reasons
upstream of any perceptual question. Chapter~\ref{ch:painting}'s
painting case illustrates the distinction cleanly: the rendering from
event to canvas is low-fidelity for causal-inference tasks by
\cref{prop:open-causal-space} alone, entirely independent of the fork; a
further rendering from canvas to a specific viewer's percept is where
the fork actually applies. By \cref{prop:admissibility-monotone},
fidelity cannot increase under further rendering, so the fork can only
leave an already-low fidelity the same or lower -- it can never resolve
an upstream ambiguity into certainty. Chess's board case
(Chapter~\ref{ch:chess}) is fork-relevant precisely because it lacks
this confound: the board's own fidelity for determining legal
continuations is high once augmented (\cref{prop:closed-sufficiency}),
so whether a player's perceptual access reaches that same fidelity is
the fork's question in comparatively pure form.
\end{remark}

\section{Worked Material Under the Narrow Reading}

The remainder of the monograph's treatment of memory
(Chapter~\ref{ch:memory}) proceeds under the narrow reading, since it is
the reading supported by existing empirical literature independent of
MEM|8's broader claim, and flags explicitly, wherever embodied perception
is discussed in Part~\ref{part:worked-domains}, that the broad reading
remains an open alternative per \cref{oq:mem8-fork}.
